\Needspace{12\baselineskip}
\begin{center}
\textit{Formen mit positiver nichtquadratischer Determinante.}
\end{center}

\Needspace{8\baselineskip}
\Article{183}
\textsc{Aufgabe.} \textit{Zu einer beliebigen Form $(a,b,a')$ mit positiver nichtquadratischer Determinante $=D$ eine eigentlich äquivalente Form $(A,B,C)$ finden, in der $B$ positiv und $<\surd D$ ist, während $A$, falls positiv, oder $-A$, falls $A$ negativ ist, zwischen $\surd D+B$ und $\surd D-B$ liegt.}

\textit{Lösung.} Wir nehmen an, dass die gegebene Form noch nicht beide Bedingungen erfüllt; sonst müsste man keine andere Form suchen. Außerdem bemerken wir, dass bei einer Form mit \textit{nichtquadratischer} Determinante der erste oder letzte Term nicht $=0$ sein kann (Artikel 171, Anmerkung). Sei $b'\equiv-b\ (\mathrm{mod.}\ a')$, und zwar zwischen den Grenzen $\surd D$ und $\surd D\mp a'$ gelegen (mit dem oberen Vorzeichen, wenn $a'$ positiv ist, und dem unteren, wenn es negativ ist). Dass dies möglich ist, beweist man leicht wie in Artikel 3. Man setzt $\frac{b'b'-D}{a'}=a''$, was eine ganze Zahl ist, weil $b'b'-D\equiv bb-D\equiv aa'\equiv0\ (\mathrm{mod.}\ a')$. Ist nun $a''<a'$, so sei wieder $b''\equiv-b'\ (\mathrm{mod.}\ a'')$, und zwar zwischen $\surd D$ und $\surd D\mp a''$ gelegen (je nachdem, ob $a''$ positiv oder negativ ist), und $\frac{b''b''-D}{{a''}}=a'''$. Ist erneut $a'''<a''$, so sei wiederum $b'''\equiv-b''\ (\mathrm{mod.}\ a''')$, zwischen $\surd D$ und $\surd D\mp a'''$ gelegen, und $\frac{b'''b'''-D}{a'''}=a''''$. Man setzt dies fort, bis die Folge $a',a'',a''',a''''$ usw. einen Term $a^{m+1}$ erreicht, der nicht kleiner als sein Vorgänger $a^m$ ist. Das muss schließlich eintreten, da man sonst eine unendliche Folge ganzer Zahlen hätte, die immer weiter abnehmen. Setzt man dann $a^m=A$, $b^m=B$, $a^{m+1}=C$, so erfüllt die Form $(A,B,C)$ alle Bedingungen.

\textit{Beweis.} I. In der Folge der Formen $(a,b,a')$, $(a',b',a'')$, $(a'',b'',a''')$ usw. ist jede ihrer Vorgängerin benachbart. Die letzte $(A,B,C)$ ist also der ersten $(a,b,a')$ eigentlich äquivalent.

II. Da $B$ zwischen $\surd D$ und $\surd D\mp A$ liegt (stets mit dem oberen Vorzeichen, wenn $A$ positiv ist, und dem unteren, wenn $A$ negativ ist), sind bei der Setzung $\surd D-B=p$, $B-(\surd D\mp A)=q$ die Zahlen $p,q$ offenbar positiv. Man bestätigt leicht $qq+2pq+2p\surd D=D+AA-BB$; also ist $D+AA-BB$ eine positive Zahl, die wir $=r$ setzen. Wegen $D=BB-AC$
\SourceLeaf{169}
ergibt sich $r=AA-AC$, also ist $AA-AC$ positiv. Da aber nach Voraussetzung $A$ nicht größer als $C$ ist, kann dies offenbar nur eintreten, wenn $AC$ negativ ist, also die Vorzeichen von $A,C$ entgegengesetzt sind. Daraus folgt $BB=D+AC<D$ und somit $B<\surd D$.

III. Da ferner $-AC=D-BB$ gilt, folgt $AC<D$ und daher (weil $A$ nicht $>C$ ist) $A<\surd D$. Also ist $\surd D\mp A$ positiv und damit auch $B$, das zwischen den Grenzen $\surd D$ und $\surd D\mp A$ liegt.

IV. Umso mehr ist dann $\surd D+B\mp A$ positiv. Da $\surd D-B\mp A=-q$ negativ ist, liegt $\pm A$ zwischen $\surd D+B$ und $\surd D-B$. \textit{W.\,z.\,b.\,w.}

\textit{Beispiel.} Die gegebene Form sei $(67,97,140)$, ihre Determinante $=29$. Hier findet man die Folge der Formen $(67,97,140)$, $(140,-97,67)$, $(67,-37,20)$, $(20,-3,-1)$, $(-1,5,4)$. Die letzte ist die gesuchte.

\Needspace{9\baselineskip}
Solche Formen $(A,B,C)$ mit positiver nichtquadratischer Determinante $D$, bei denen $A$, positiv genommen, zwischen $\surd D+B$ und $\surd D-B$ liegt, während $B$ positiv und $<\surd D$ ist, nennen wir \textit{reduzierte Formen}. Die reduzierten Formen mit positiver nichtquadratischer Determinante unterscheiden sich also etwas von den reduzierten Formen mit negativer Determinante. Wegen der großen Analogie zwischen beiden wollten wir aber keine verschiedenen Namen einführen.

\Needspace{8\baselineskip}
\Article{184}
Könnte man die Äquivalenz zweier \textit{reduzierter} Formen mit positiver Determinante ebenso leicht erkennen wie bei Formen mit negativer Determinante (Artikel 172), so ließe sich ohne Mühe entscheiden, ob zwei \textit{beliebige} Formen mit derselben positiven Determinante äquivalent sind. Hier ist die Lage aber ganz anders: Sehr viele reduzierte Formen können zueinander äquivalent sein. Bevor wir diese Aufgabe angehen, müssen wir die Natur der reduzierten Formen genauer untersuchen (mit positiver nichtquadratischer Determinante, was hier immer mitgemeint ist).

1) Ist $(a,b,c)$ eine reduzierte Form, so haben $a$ und $c$ entgegengesetzte Vorzeichen. Setzt man nämlich die Determinante der Form $=D$, so gilt $ac=bb-D$, was wegen $b<\surd D$ negativ ist.

2) Die Zahl $c$ liegt ebenso wie $a$, positiv genommen, zwischen $\surd D+b$ und $\surd D-b$. Denn $-c=\frac{D-bb}{a}$; also liegt $c$, wenn man vom Vorzeichen absieht, zwischen $\frac{D-bb}{\surd D+b}$ und $\frac{D-bb}{\surd D-b}$, \textit{das heißt} zwischen $\surd D-b$ und $\surd D+b$.

3) Damit ist klar, dass auch $(c,b,a)$ eine reduzierte Form ist.

\SourceLeaf{170}
4) Sowohl $a$ als auch $c$ sind $<2\surd D$. Beide sind nämlich $<\surd D+b$ und damit umso mehr $<2\surd D$.

5) Die Zahl $b$ liegt zwischen $\surd D$ und $\surd D\mp a$ (mit dem oberen Vorzeichen, wenn $a$ positiv ist, und dem unteren, wenn es negativ ist). Da nämlich $\pm a$ zwischen $\surd D+b$ und $\surd D-b$ liegt, ist $\pm a-(\surd D-b)$ oder $b-(\surd D\mp a)$ positiv; $b-\surd D$ dagegen ist negativ. Also liegt $b$ zwischen $\surd D$ und $\surd D\mp a$. --- Genau auf dieselbe Art beweist man, dass $b$ zwischen $\surd D$ und $\surd D\mp c$ liegt (je nachdem, ob $c$ positiv oder negativ ist).

6) \textit{Jede reduzierte Form $(a,b,c)$ hat auf jeder Seite genau eine reduzierte Nachbarform.}

Man setzt $a'=c$, $b'\equiv-b\ (\mathrm{mod.}\ a')$, zwischen $\surd D$ und $\surd D\mp a'$ gelegen\HistoricalNote{*)}{Wo Vorzeichen zur Wahl stehen, gilt stets das obere, wenn $a'$ positiv ist, und das untere, wenn $a'$ negativ ist.}, und $c'=\frac{b'b'-D}{a'}$. Dann ist die Form $(a',b',c')$ der Form $(a,b,c)$ nach rechts benachbart. Zugleich ist klar, dass eine reduzierte Form, die der Form $(a,b,c)$ nach rechts benachbart ist, keine andere als $(a',b',c')$ sein kann. Dass sie tatsächlich reduziert ist, beweisen wir so:

A) Setzt man
\[
\surd D+b\mp a'=p,\quad\pm a'-(\surd D-b)=q,\quad\surd D-b=r
\]
so sind diese $p,q,r$ nach (2) oben und nach der Definition einer reduzierten Form positiv. Weiter setzt man
\[
b'-(\surd D\mp a)=q',\quad\surd D-b'=r'
\]
\Needspace{10\baselineskip}
Dann sind $q',r'$ positiv, weil $b'$ zwischen $\surd D$ und $\surd D\mp a'$ liegt. Schließlich sei $b+b'=\pm ma$; dann ist $m$ eine ganze Zahl. Offenbar gilt $p+q'=b+b'$, also ist $b+b'$ oder $\pm ma'$ positiv und damit auch $m$. Daraus folgt, dass $m-1$ sicher nicht negativ ist. Weiter ergibt sich
\[
r+q'\pm ma'=2b'\pm a',\quad\text{oder}\quad2b'=r+q'\pm(m-1)a'
\]
also sind $2b'$ und $b'$ notwendig positiv. Da $b'+r'=\surd D$, gilt $b'<\surd D$.

B) Weiter ergibt sich
\[
r\pm ma'=\surd D+b',\quad\text{oder}\quad r\pm(m-1)a'=\surd D+b'\mp a'
\]
also ist $\surd D+b'\mp a'$ positiv. Daraus und aus $\pm a'-(\surd D-b')=q'$,
\SourceLeaf{171}
was also positiv ist, folgt, dass $\pm a'$ zwischen $\surd D+b'$ und $\surd D-b'$ liegt. --- Somit ist $(a',b',c')$ eine reduzierte Form.

Ebenso beweist man: Setzt man ${}'c=a$, ${}'b\equiv-b\ (\mathrm{mod.}\ {}'c)$, zwischen $\surd D$ und $\surd D\pm{}'c$ gelegen, und ${}'a=\frac{{}'b{}'b-D}{{}'c}$, so ist die Form $({}'a,{}'b,{}'c)$ reduziert. Offenbar ist diese Form der Form $(a,b,c)$ nach links benachbart, und keine andere reduzierte Form als $({}'a,{}'b,{}'c)$ kann diese Eigenschaft haben.

\textit{Beispiel.} Die reduzierte Form $(5,11,-14)$ mit Determinante $=191$ hat nach rechts die reduzierte Nachbarform $(-14,3,13)$ und nach links diese $(-22,9,5)$.

7) Ist der reduzierten Form $(a,b,c)$ die reduzierte Form $(a',b',c')$ nach rechts benachbart, so ist der reduzierten Form $(c,b,a)$ die Form $(c',b',a')$ nach links benachbart. Ist der reduzierten Form $(a,b,c)$ die Form $({}'a,{}'b,{}'c)$ nach links benachbart, so ist der reduzierten Form $(c,b,a)$ die reduzierte Form $({}'c,{}'b,{}'a)$ nach rechts benachbart. Außerdem sind die Formen $(-{}'a,{}'b,-{}'c)$, $(-a,b,-c)$, $(-a',b',-c')$ reduziert, und die zweite ist der ersten, die dritte der zweiten nach rechts benachbart; oder die erste ist der zweiten und die zweite der dritten nach links benachbart. Ebenso verhält es sich mit den drei Formen $(-c',b',-a')$, $(-c,b,-a)$, $(-{}'c,{}'b,-{}'a)$. Dies ist so klar, dass es keiner weiteren Erläuterung bedarf.
